% ---
% title: 정리 틀
% description: 번호 붙는 정리 환경 (한글 이름)
% place: last
% ---
% 연구 전용 서식 뒤에 불러오며, 아직 없는 환경만 만든다.
\makeatletter
\newcommand{\rwtheorem}[3]{\@ifundefined{#1}{\theoremstyle{#2}\newtheorem{#1}{#3}[section]}{}}
\makeatother
\rwtheorem{definition}{definition}{정의}
\rwtheorem{theorem}{plain}{정리}
\rwtheorem{lemma}{plain}{보조정리}
\rwtheorem{proposition}{plain}{명제}
\rwtheorem{corollary}{plain}{따름정리}
\rwtheorem{remark}{remark}{주의}
